\summarychapter{48}{Windows to the Christian World}{Places and Persons of Faith}

\begin{inbrief}
The last chapter adds texture to the general account of the Global South by telling the stories of a
few places and people: the survival and growth of the church in China, the steadfastness of Cardinal
Kung, the rise of Korean Christianity and its spirituality of prayer, and the East African Revival.
In each, a church indebted to missionaries has become richly its own, has grown remarkably and has
produced believers who bear suffering and extend grace.
\question{What do the stories of particular places and persons show about the church of the Global
South?}
\end{inbrief}

The chapter opens with a caution. It is humbling to remember that the Global Church is being
described here in Western terms, and, as Randy Richards has written, ``generalizations are always
wrong and usually helpful.'' Some texture can be added by meeting, however briefly, a few of the
places and people who together make up the people of God.

\ssection{China}
China, treated first and at greatest length, had known great violence over religion; one religious
uprising cost twenty-five million lives. Christian missions were received with suspicion, and
missionaries were expelled in the ninth and thirteenth centuries. The Jesuits had more lasting but
modest success in the 1600s, and other Christians complained that they had compromised Christian
teaching to make it accessible to the Chinese. Western nations made China the focus of their missions
in the late 1800s and early 1900s. After the Second World War civil war broke out between the
Nationalists under Chiang Kai-shek and the Communists under Mao Zedong. The Nationalists were driven
from the mainland to Taiwan, and the Communists formed the People's Republic and in 1949 expelled the
missionaries. Many wondered whether Christians could survive, given China's past, the suspicion that
Christianity was a tool of Western imperialism, the weak and nominal state of some converts and the
Communists' hostility to religion.

At first the aim was control. Mao believed that faith would die a natural death once the oppression
and manipulation practiced by religion were restrained. Governing bodies were created to ensure
patriotic compliance: the Protestant Three-Self Patriotic Movement and the Catholic Patriotic
Association (dated 1955 and 1957 in the book) became powerful agencies that censured and coerced in
the name of loyalty to the nation. Chinese Christians were badly divided. Some joined and bore the
difficulties so that a Christian voice might survive, and others protested that such cooperation
compromised the faith. The Protestant body took its name from an old missionary strategy for
self-reliant churches, ``self-governing, self-supporting, and self-propagating,'' words that had a
further meaning for Communists, who saw all Western influence as oppressive. Those who cooperated
were regarded as traitors; theirs was the ``three-self church,'' now more often called the
registered church. The others risked life and limb for their independence: the underground or house
church, now usually called the unregistered church. The government persecuted brutally and also
counted on isolation from the West to bring Chinese Christianity down. A later policy tried to end
religion by force. Mao himself set off the Great Proletarian Cultural Revolution (1966--1976), with
its massive brutality, re-education and cult-like veneration of the chairman. Church property was
confiscated, and even Christians of the registered churches were persecuted.

After Mao's death more moderate voices gained influence, embodied in Deng Xiaoping. President Jimmy
Carter normalized relations with China in January 1979. Visiting the United States, Deng thanked him
and asked whether he could do Carter a kindness in return. Carter said that as a boy he had given
five cents a week to build hospitals and schools in China and had thought missionaries the heroes of
his faith; he noted that China forbade missionaries, Bibles and freedom of worship and asked that
all three be allowed. Deng asked for time and answered the next morning. China would not admit
missionaries, but it would authorize the distribution of Bibles and pass a law allowing free
worship. The printing and distribution of Bibles came quickly, and in 1982 a law began the process
of giving Christians greater freedom to worship. Actual freedom varies widely by region, but more
and more unregistered churches meet openly and registered churches enjoy greater liberty. Hostility
between the two is diminishing and cooperation is becoming more usual. A third influence has
appeared as well. China's economic success has drawn entrepreneurs from around the world, and
foreigners with passports have the legal right to worship and assemble freely; their international
churches are an important spiritual influence and a link to the other churches. The number of
converts is mind-bending. A conservative estimate for 2000 is 89 million believers, 7.1 percent of
the population.

\ssection{Cardinal Kung: Faithful Resolve}
Ignatius Kung Pin-Mei (also Gung Pin-Mei) was ordained in 1930 and in 1950 became the first
native-born Chinese bishop of the great port of Shanghai. He was imprisoned in 1955 for refusing to
renounce the pope and to accept Chinese control of the Catholic Church. The authorities paraded him
before a huge crowd at the dog-racing stadium to denounce the pope, and he shouted instead, ``Long
live Christ the King, long live the pope.'' After five years in prison he was convicted of treason in
1960 and for the next twenty-five years was kept in isolation except at times of physical labor,
without Bible or rosary, letters or contact with his family.

As China moved slowly toward a more open economy and, more slowly still, a more open culture, Deng
Xiaoping arranged a banquet in 1984 for Cardinal Jaime Sin of Manila, the Filipino churchman known
for his part in the popular uprising that unseated Ferdinand Marcos. As a gesture the old bishop,
almost thirty years a prisoner, was invited. The two were allowed no private moment, but Kung seized
the chance to offer a song and defiantly sang in Latin the words of Matthew 16:18, ``Upon the rock I
will build my church.'' His resolve had survived thirty years of labor and isolation. The next year
he was released for ``health reasons.'' Not even Cardinal Sin knew that evening that Pope John Paul
II had already named Kung a cardinal \emph{in pectore}, ``in the heart,'' a secret appointment made
when publicity would be inappropriate or dangerous. Kung had first tried to work with the
authorities and later defied them. Many worried that the Jesuits had compromised Christianity, and
others wondered how any true believer could work with the registered church; yet time has shown the
Spirit fashioning strong faith from humble beginnings.

\ssection{Korea}
Korea's story is unusual, for with some qualification Koreans brought Christianity to Korea. In 1784
Yi Sung-hun was baptized while on a diplomatic mission to China and came home with writings of the
Jesuit missionary Matteo Ricci. Having no priest, the Korean Christians baptized one another and
held Mass. In 1794 a priest from China, James Chou, was sent in secretly and ministered to almost
four thousand practicing Christians. These Catholics were suspected of ties with foreign powers, and
waves of martyrdom, in which Chou himself died, continued for nearly half a century. In 1873 King
Kojong showed greater openness to Western influence, and Catholics increased. John Ross, a Scottish
Presbyterian missionary in China, translated the New Testament into Korean, and from the 1880s many
Protestant missionaries began work.

Their work reached a young man named Kil Sun-ju (Sun Chu Kil in the book). His family moved to
Pyongyang in 1885, the year several American Presbyterian missionaries arrived. They supported Bible
translation and followed the strategy of John Nevius, an American missionary in China, for strong,
self-reliant churches. Their coming coincided with Korean disillusionment at living in the shadow of
China, Russia and Japan. Kil had struggled with his Taoism and begun to explore Christianity. The
missionary Samuel Moffett invited him to read \emph{Pilgrim's Progress}, and after a night of
weeping and prayer he was converted. He became active in church leadership and in evangelism among
working people and in 1903 entered the Pyongyang Presbyterian Seminary, which Moffett supported.
Reports of the Welsh revival moved him to hold morning prayer meetings for a revival in Korea. Near
the end of a ten-day training series for new converts a missionary publicly confessed ``stubbornness
and pride.'' Many others followed, and Kil himself admitted mishandling a will. The meeting lasted
from eight in the evening until two in the morning, with the Koreans praying simultaneously, each
aloud at the same time, a symphony of voices that some say is pleasing to God. Westerners have often
noted this as distinctive of Korean spirituality, though it is found elsewhere.

Japan annexed Korea in 1910 and ruled it until the end of the Second World War, treating the
churches brutally. Many Christians led in the call for independence, and Kil was a powerful voice in
the movement while holding to Christian nonviolence. For once Christians earned a name for standing
against colonization. He was appointed to an influential pastorate, Central Presbyterian Church. His
part in the independence movement finally cost his life and his son's. When his fellow leaders were
arrested he gave himself up, and after his release he died of the effects of his mistreatment. Mark
Noll reports that he preached more than seventeen thousand sermons and shared in founding sixty
churches.

After the war Korea was divided between a North backed by the Soviet Union and a South backed by the
United States. Fighting ceased in 1953, but the Communists struck at the church, and two-thirds of
Korean Christians were then in the North. The secrecy of its government makes the state of
Christianity there very hard to assess. In the South Presbyterians, Pentecostals and more recently
Catholics have grown remarkably, the Korean revival acting as a catalyst. In 1914 one Korean in a
hundred was Protestant, and by 2010 one in three was Christian. Korea has the world's largest
churches. The Yoido Full Gospel Church has well over a million members, and its cell structure is a
model for churches around the world. South Korea is second only to the United States in founding
seminaries and sending missionaries, who carry with them the policy of self-reliant churches and a
distinctive spirituality rooted in prayer. Visitors remark on the simultaneous prayer and on the
Prayer Mountain, where Koreans book time at retreat centers for prayer and fasting. David Yonggi
Cho, pastor of the Yoido church, credits its growth not to its cell structure or its efficiency but
to its prayer.

\ssection{East African Revival}
One part of the great revival that swept East Africa is the friendship of two men who had both known
discouragement. Simeon Nsibambi, born in 1897 and converted in 1922, had reason to hope that he
would be chosen to study abroad, but it was not to be. God met his severe disappointment in a vision
that asked him to weigh the scholarship against the gospel and the forgiveness he had received.
Joseph Church was converted as a Cambridge student in August 1920. Cambridge had known evangelical
stirrings. D.~L.~Moody had preached there almost forty years earlier, and afterward many students
were called to missions, among them the Cambridge Seven, who went to China, and George Lawrence
Pilkington, who witnessed in Uganda in 1893 a revival that foreshadowed the later one. Church worked
in Africa as a doctor and proved an able missionary.

By 1929 he was deeply discouraged in body and spirit and longing for companionship when he met
Nsibambi. They shared their longing for revival in themselves and in Africa and for two days prayed
and studied the Bible together, an experience that changed both lives. Church remembered it as the
time when God ``brought me to the end of myself and thought fit to give me a share in the power of
Pentecost.'' The fellowship widened again and again. In September 1935 another group met for
renewal, shaped by the Keswick theme that complete surrender to Christ comes before spiritual
breakthrough, and planned a series of meetings that may be counted the starting point of the
revival. It included Church, the only white member, and Nsibambi; William Nagenda soon joined and
did more than anyone to spread the revival. Two words show its character: \emph{abaka}, ``on fire,''
and \emph{balokole}, ``the saved.''

The core leaders are credited not only with preaching but with the organizational wisdom of sending
out teams that included lay people. The revival moved from Uganda into Kenya, Rwanda, Burundi and
Tanzania (then Tanganyika) and touched many beyond the region. The numbers defy imagination. Near
the end of the century the Christian share of the population was 89 percent in Uganda, 79 in Kenya,
81 in Rwanda, 90 in Burundi and 51 in Tanzania. Nagenda and Church were perhaps the most effective
team, and Nagenda's strong family vindicated his powerful preaching. He was also part of an
important test. In 1941 twenty-nine students were expelled from Bishop Tucker Theological College by
leaders alarmed at the zeal of these \emph{balokole}. The school's authorities, all white, were able
to restore almost every one, Nagenda included, so that students could express their fervor and stay
loyal to their denomination, a flexibility that proved very important. The movement also crossed
denominational lines: Brethren, Seventh-day Adventists, Mennonites, Baptists and Methodists all
report taking part.

Another celebrated voice was Festo Kivengere. Born into a ruling family, he worked as a shepherd boy
and read children's stories about Jesus as he watched the flock. He was converted at an evangelistic
meeting, studied theology, served as a pastor and became an Anglican bishop. He was a friend of
Billy Graham, translated for him and appeared at his crusades. His courage before Idi Amin was
widely admired. Amin murdered his archbishop, Janani Luwum, and Kivengere stayed in exile until Amin
fell, becoming in those years an internationally recognized voice for social justice and for all of
Africa. His theology was firmly rooted in the cross, and he believed that Christ called him to love
Amin despite his atrocities.

China, Korea and East Africa have several things in common. Each has a rich array of Christian
cultures of its own while remaining indebted to missionary efforts. Each has seen remarkable
numerical growth. And each has produced believers who give themselves to Christ in bearing suffering
and extending grace.

\begin{keyfacts}
\keypeople{Mao Zedong; Deng Xiaoping; Jimmy Carter; Cardinal Ignatius Kung; Cardinal Jaime Sin; Yi
Sung-hun; James Chou; John Ross; Samuel Moffett; Kil Sun-ju; David Yonggi Cho; Simeon Nsibambi;
Joseph Church; William Nagenda; Festo Kivengere; Janani Luwum.}
\keydates{%
  \kd{1784}{Yi Sung-hun baptized in China; Christianity enters Korea through Koreans}
  \kd{1885}{American Presbyterian missionaries arrive in Pyongyang}
  \kd{1910}{Japan annexes Korea}
  \kd{1929}{Joseph Church and Simeon Nsibambi meet}
  \kd{1935}{Meetings that begin the East African Revival}
  \kd{1949}{People's Republic of China expels the missionaries}
  \kd{1955}{Bishop Kung of Shanghai imprisoned}
  \kd{1966--1976}{Cultural Revolution in China}
  \kd{1979}{The United States normalizes relations with China; Bibles and freedom of worship
    promised}
  \kd{1985}{Kung released after thirty years}}
\keyterms{Registered and unregistered churches; three-self principle; house church; Cultural
Revolution; \emph{in pectore}; self-reliant churches; simultaneous prayer; Prayer Mountain; cell
structure; East African Revival; \emph{balokole}.}
\end{keyfacts}
